\chapter{Conclusion}

This diploma project addressed the problem of automating external penetration testing of web applications. The work proposed and implemented \textbf{ReconX} --- a method and a Python~3 framework that integrates reconnaissance, configuration analysis, and a broad set of OWASP-aligned active probes into a single reproducible pipeline, normalizes their output into a unified finding schema, and uses a local LLM to produce an analyst-readable report.

\textbf{Results acquired.}
\begin{itemize}
  \item The theoretical part of the work analyzed the OWASP Top~10 (2021) catalogue, the structure of a modern external-pentest engagement, the main classes of web vulnerabilities, and the open-source toolchain used by the offensive-security community.
  \item Functional and non-functional requirements for an integrated external-pentest framework were formulated, including reproducibility, modularity, scope enforcement, finding normalization, and LLM-assisted reporting.
  \item A modular pipeline architecture was designed, consisting of 46 specialized modules grouped into 15 sequential execution stages, with intra-stage thread-level parallelism.
  \item A central \textit{finding registry} was implemented that maps every probe's raw output into a unified record with severity, confidence, CWE, CVSS, OWASP, and MITRE ATT\&CK metadata, and that derives a triage \textit{verdict} of \textit{confirmed} versus \textit{candidate}.
  \item Reconnaissance modules were implemented: subdomain discovery from nine passive sources, DNS auditing, port scanning, technology and CMS fingerprinting, endpoint and parameter discovery, source-map analysis, and OpenAPI parsing.
  \item Configuration-analysis modules were implemented for TLS, security headers, CORS, authentication cookies, JWT tokens, OpenAPI security anti-patterns, and secret hunting in public GitHub repositories.
  \item Active probes were implemented for reflected XSS, SQL injection (via sqlmap), IDOR with cross-profile comparison, host header injection, cache poisoning, prototype pollution, open redirect, SSRF / SSTI, XXE, HTTP request smuggling, WebSocket security, race conditions, GraphQL auditing, and OAuth misconfigurations.
  \item A correlator implementing nine cross-finding rules and an asset-risk module producing a ranked per-host risk score were implemented, allowing the framework to highlight the most exploitable assets at the top of the report.
  \item A local LLM integration (Ollama with DeepSeek-R1~7B as the default model and an OpenAI-compatible fallback) was implemented and produces an executive-summary report from the structured findings.
  \item An HTML report renderer, a SIEM-compatible JSON export, a per-session append-only audit log of every HTTP request, and a Telegram notifier for long-running scans were implemented.
  \item The framework was experimentally evaluated in two complementary settings: a controlled laboratory run against \textit{OWASP Juice Shop} and an authorized real-world case study against the \texttt{miras.app} academic management system. The detected vulnerabilities matched the expected categories of both targets, the scope enforcer correctly blocked the single out-of-scope request observed during the real scan, the asset-risk module placed the most exposed administrative subdomain at the top of the ranking, and the full real-target pipeline completed in approximately 1~minute~43~seconds.
\end{itemize}

\textbf{Improvements that can be made in future work.}
\begin{itemize}
  \item Add a headless-browser-based DAST stage that drives the application through Playwright in addition to the current HTTP-level probing, in order to discover client-side issues that require an actual JavaScript runtime.
  \item Extend the active-probe set to cover NoSQL injection, GraphQL batch-based authorization bypass, mass-assignment, and BOLA on REST APIs with finer granularity.
  \item Introduce a feedback loop in which the LLM proposes follow-up probes (for example, targeted SSRF callbacks against an interactsh server) that are then automatically dispatched by the framework, instead of leaving the proposal to the human analyst.
  \item Replace the current per-rule correlator with a learned ranker trained on labeled finding sets, in order to improve the prioritization of multi-signal observations.
  \item Build a dedicated web dashboard --- comparable to OWASP DefectDojo --- on top of the SIEM-compatible JSON export, with multi-scan diffing and historical trend analysis.
  \item Package the framework as an installable distribution and add a containerized runtime so that the framework can be deployed both on a local Kali-Linux workstation and inside a CI/CD pipeline.
\end{itemize}

The work demonstrates that external penetration testing of web applications can be substantially automated without sacrificing explainability, and that local LLMs are mature enough to be integrated as a first-class component of an offensive-security pipeline.
